\begin{table}[H]
\centering
\caption{Variación Regional en las Defecciones de la Coalición}
\label{tab:region}
\small
\begin{adjustbox}{max width=\textwidth}
\begin{tabular}{lS[table-format=4]S[table-format=2.1]S[table-format=2.1]S[table-format=2.1]S[table-format=2.1]S[table-format=2.2]S[table-format=2.1]S[table-format=2.1]}
\toprule
{\textbf{Región}} & {\textbf{N}} & \multicolumn{2}{c}{\textbf{CA}} & \multicolumn{2}{c}{\textbf{PC}} & \multicolumn{2}{c}{\textbf{PN}} & {\textbf{Blancos}} \\
\cmidrule(lr){3-4} \cmidrule(lr){5-6} \cmidrule(lr){7-8}
 & & {\textbf{$\rightarrow$FA (\%)}} & {\textbf{$\rightarrow$PN (\%)}} & {\textbf{$\rightarrow$FA (\%)}} & {\textbf{$\rightarrow$PN (\%)}} & {\textbf{$\rightarrow$FA (\%)}} & {\textbf{$\rightarrow$PN (\%)}} & {\textbf{(\%)}} \\
\midrule
Área Metropolitana & 3671 & 47.4 & 44.5 & 7.0 & 91.1 & 0.17 & 98.1 & 8.2 \\
 & & {[41.8, 52.8]} & {[39.1, 50.0]} & {[5.5, 8.6]} & {[89.5, 92.8]} & {[0.004, 0.59]} & {[97.6, 98.6]} & {[6.5, 9.8]} \\
Interior & 3600 & 55.5 & 41.0 & 10.7 & 87.1 & 8.1 & 90.1 & 3.5 \\
 & & {[51.0, 59.9]} & {[36.4, 45.5]} & {[9.6, 11.7]} & {[86.0, 88.2]} & {[7.5, 8.7]} & {[89.5, 90.7]} & {[2.2, 4.9]} \\
\bottomrule
\end{tabular}
\end{adjustbox}
\begin{flushleft}
\small \textit{Nota:} CA, Cabildo Abierto; PC, Partido Colorado; PN, Partido Nacional; FA, Frente Amplio. Área Metropolitana incluye Montevideo y Canelones. Los valores representan medias posteriores con intervalos creíbles al 95\% entre corchetes.
\end{flushleft}
\end{table}
